\section{问题分析与统一建模框架}
\subsection{从运载能力到可执行方案}
山区应急投送的关键困难，是一条地理上可达的航线未必能在给定电量、交付时限和通信条件下执行。题目提供一个调度中心、\NumServices{} 个服务区与 \NumBoxes{} 个不可拆分货箱，要求从单点往返逐步扩展至异构运输调度、空中中继和独立任务分区\cite{problem}。本文以官方标准测试场景为研究对象，不把背景灾情叙述直接作为需求预测数据。

运输路径与能量消耗不能分开处理。既有无人机配送研究已指出，载荷相关能耗和多架次复用会影响路线可行性\cite{dorling}；不同能耗模型的使用条件也不能互相替代\cite{zhang}。本题的重点不是重新拟合动力学，而是在官方统一规则下，让每个阶段的载荷、高度、用电、充电和通信判定保持一致。

我们把四问组织为一条逐步增加约束的计算链。Q1 给出不受实体资源调度影响的能力基线；Q2 加入无人机、电池和交付时限，得到可执行的运输候选；Q3 在固定离散候选上联动运输开始时刻与中继服务窗口；Q4 保持 Q3 任务不变，度量跨组资源共享被取消的代价。四问的可行域不同，因此架次数和能耗的差异不能直接解释为同一优化问题的算法优劣。

\subsection{统一思想与证据层级}
统一框架如图\ref{fig:framework}。底层把每个航段转化为可复用的时间、能量和轨迹对象，上层只决定货箱、路径、资源与时序。相同的物理对象再交给验证器重算，避免求解模块自己报告成功后直接进入论文。
\paperfig{fig01_framework}{同一物理计算连接能力边界与资源、通信时序。左侧为真实 DEM 剖面和安全载荷曲线，右侧为实际运输与中继资源占用；中继浅色条表示飞行与准备，深色条表示通信服务。Q1 为精确组批，Q2 为启发式候选，Q3 的线性规划最优性仅针对固定离散方案。}{framework}

\FloatBarrier
